\chapter{Conclusiones}\label{cap:7}

\pendiente{capítulo nuevo. Responder cada objetivo específico en un párrafo.
Viñetas de partida:}

\begin{itemize}
\tightlist
\item
  \textbf{Objetivo 1:} CR2MET muestra un gradiente norte--sur marcado: rachas medias de 65,6 días en Loa frente a 6,2 días en Biobío, con colas de cientos de días en el norte.
\item
  \textbf{Objetivo 2:} igualar la frecuencia de días húmedos es necesario, pero no suficiente: ALADIN reproduce el patrón espacial, pero acorta las rachas, sobre todo en la cola (Δ media −8,2 días; Δ t99 −46,3 días).
\item
  \textbf{Objetivo 3:} hacia 2040--2074, ALADIN proyecta menos rachas, pero más largas; el número de rachas de al menos 20 días aumenta en las cuatro regiones (significativo en Coquimbo y La Araucanía).
\item
  \textbf{Objetivo 4:} la evaluación histórica es robusta al regrillado, y la dirección de la señal futura se mantiene con otro umbral y con SSP3-7.0.
\item
  \textbf{Mensaje general:} el secado proyectado es una reorganización temporal hacia períodos secos más persistentes, que no se deduce del cambio en la precipitación total.
\item
  \textbf{Trabajo futuro:} otros RCM del dominio CHP, horizontes de fin de siglo, mecanismos de circulación e impactos hidrológicos.
\end{itemize}
